\subsection{实值函数}

定义 1. 一个集 X 到实直线中的映射称为定义在 X 上的实值函数（或实函数）。

通过与 § 4 第 2 目所述类似的用语习惯，X 到 \(\overline{\mathbf{R}}\) 中的映射在本节及下一节中也称为定义在 X 上的实值函数。X 到 \(\mathbf{R}\) 中的映射将称为有限实值函数。

若 \(f\) 与 \(g\) 是定义在 \(\mathbf{X}\) 上的两个实值函数，则关系 \(f \leqslant g\) 依定义等价于“\(f(x) \leqslant g(x)\) 对所有 \(x \in \mathbf{X}\) 成立；”这个关系是集 \(\overline{\mathbf{R}}^{\mathbf{X}}\) （ \(\mathbf{X}\) 上所有实值函数的）上的一个序。由这个关系排序的集 \(\overline{\mathbf{R}}^{\mathbf{X}}\) 是一个格；因为若 \(f\) 与 \(g\) 是任意两个实值函数，则函数 \(h\) 由 \(h(x) = \sup (f(x), g(x))\) （对所有 \(x \in \mathbf{X}\) ）所定义，它是 \(\mathbf{X}\) 上既 \(\geqslant f\) 又 \(\geqslant g\) 的实值函数中的最小者；按照一般记号，我们把这个函数（即 \(f\) 与 \(g\) 在 \(\overline{\mathbf{R}}^{\mathbf{X}}\) 中的上确界）记作 \(\sup (f, g)\) 。类似地，在每个 \(x \in \mathbf{X}\) 处取值 \(\inf (f(x), g(x))\) 的函数记作 \(\inf (f, g)\) 。

注意 sup \((f, g)\) 是映射

\[
(u, v) \mapsto \sup (u, v)
\]

\(\overline{\mathbf{R}}\times \overline{\mathbf{R}}\) 到 \(\overline{\mathbf{R}}\) 中的映射与映射 \(x\to (f(x),g(x))\) （ \(\mathbf{X}\) 到 \(\overline{\mathbf{R}}\times \overline{\mathbf{R}}\) 中的）的复合。inf \((f,g)\) 亦然。

实值函数 \(f\) （定义在集 \(\mathbf{X}\) 上的）称为在 \(\mathbf{X}\) 中上有界（相应地下有界）的，如果 \(f(\mathbf{X})\) 是 \(\mathbf{A}'' = [-\infty, +\infty[\) 的子集且上有界（相应地，如果 \(f(\mathbf{X})\) 是 \(\mathbf{A}' = ] - \infty, +\infty]\) 的子集且下有界）。\(f\) 称为在 \(\mathbf{X}\) 中有界的，如果它既上有界又下有界，也就是说 \(f(\mathbf{X})\) 是 \(\mathbf{R}\) 的有界子集。

因此每个有界函数都是有限的。其逆不真，例如函数 \(\mathbf{1} / x\) （在 \(\mathbb{R}_+^* = ]0, +\infty[\) 上）所示。
% semantic-fragments: legacy-input-seam
\relax{}
